\documentclass[11pt]{article}

\usepackage[a4paper,margin=1in]{geometry}
\usepackage{amsmath,amssymb}
\usepackage{booktabs}
\usepackage{enumitem}
\usepackage{hyperref}
\usepackage{graphicx}
\usepackage{xcolor}

\title{A Tri-System Neurosymbolic AGI Prototype:\\Design Draft and Experimental Program}
\author{Project Draft (Living Document)}
\date{\today}

\begin{document}
\maketitle

\begin{abstract}
This document is a living design draft for a tri-system neurosymbolic AGI prototype.
It summarizes architecture, working-memory semantics, rule forms, learning objectives,
and an experimental plan focused on compute-efficient common-sense training.
The draft is intended for both communication and internal memory, and should be updated
as algorithms and results evolve.
\end{abstract}

\section{Motivation and Scope}
The project goal is not narrow benchmark hacking, but a scalable and interpretable learning architecture
for AGI-like reasoning under practical compute constraints.
Arithmetic benchmarks (e.g., GSM-style tasks) are treated as diagnostic slices, not the sole objective.

\subsection{Current strategy}
\begin{itemize}[leftmargin=1.5em]
  \item Keep the symbolic substrate explicit and inspectable.
  \item Allow broad rule exploration with soft, relative scoring.
  \item Use common-sense and semantic grounding as primary capability drivers.
  \item Evaluate by capability \emph{and} cost efficiency (performance per compute dollar).
\end{itemize}

\section{System Overview}
The architecture has three interacting systems sharing a common working memory (WM) and long-term memory (LTM).

\subsection{System 1: Logic Transformer (LT)}
\begin{itemize}[leftmargin=1.5em]
  \item Neural proposal and pattern-learning subsystem.
  \item Produces soft hypotheses and candidate WMEs/rules with confidence/activation.
  \item Supports autoregressive (AR) objectives and RL-driven policy/credit updates.
\end{itemize}

\subsection{System 2: Forward-Chaining Rule Engine}
\begin{itemize}[leftmargin=1.5em]
  \item Deterministic matching and rule firing over WMEs (Rete-style agenda semantics).
  \item Performs consistency-oriented symbolic propagation and truth-maintenance routines.
  \item Works naturally with heterogeneous predicate arity.
\end{itemize}

\subsection{System 3: Backward-Chaining / Abductive Layer}
\begin{itemize}[leftmargin=1.5em]
  \item Goal-directed reasoning and hypothesis testing.
  \item Supports belief revision logic and abductive completion.
  \item Feeds revised beliefs/rules back to WM and System 2.
\end{itemize}

\section{Working Memory and Knowledge Representation}

\subsection{WME form}
A Working Memory Element (WME) is represented as:
\[
\texttt{WME} = (\texttt{concept}, \texttt{args/fields}, \texttt{activation}, \texttt{source}, \texttt{layer}, \texttt{ttl}, \texttt{provenance})
\]
Predicate arity is \textbf{not globally uniform}: different predicates may have different numbers of arguments.
The only arity constraint is per-predicate consistency (e.g., \texttt{has/3} remains arity-3 across uses).

\subsection{Heterogeneous rule/fact shapes}
From symbolic logic and engine semantics, heterogeneous arity is acceptable.
For tensor training efficiency in LT, shape management can be done by:
\begin{enumerate}[leftmargin=1.5em]
  \item shape bucketing (group by template shape), or
  \item padding/masking within a bucket.
\end{enumerate}

\section{Rule Language and Recommended Shapes}

\subsection{Canonical rule schema}
\[
\texttt{rule} := (\texttt{head}, \texttt{body}, \texttt{calc}, \texttt{target}, \texttt{metadata})
\]
with
\begin{itemize}[leftmargin=1.5em]
  \item \texttt{head}: predicate and argument list,
  \item \texttt{body}: conjunction of predicates and argument lists,
  \item \texttt{calc}: optional arithmetic/functional expression,
  \item \texttt{target}: bound output variable.
\end{itemize}

\subsection{Practical shape library for AGI training}
To balance expressivity and batching efficiency, use a small template library:
\begin{enumerate}[leftmargin=1.5em]
  \item Unary type/state: \texttt{type(X,T)}, \texttt{flag(X,S)}
  \item Binary relations: \texttt{isa(X,Y)}, \texttt{part\_of(X,Y)}, \texttt{causes(X,Y)}
  \item Ternary relations: \texttt{rel(S,R,O)}, \texttt{attr(X,A,V)}
  \item Event-role reification:
  \begin{itemize}
    \item \texttt{event(E,Kind)}
    \item \texttt{role(E,Role,Arg)}
    \item \texttt{time(E,T)}, \texttt{loc(E,L)}
  \end{itemize}
  \item Quantity/change forms:
  \begin{itemize}
    \item \texttt{qty(Owner,Item,Value,Unit)}
    \item \texttt{delta(Event,Owner,Item,DV)}
    \item \texttt{state(Owner,Item,Value,Time)}
  \end{itemize}
\end{enumerate}

\subsection{Design note}
Mixed-size rule sets are acceptable and often desirable.
For GPU throughput, prefer a \emph{small number of recurrent shapes} rather than unconstrained free-form arity explosion.

\section{Learning Objectives}

\subsection{Autoregressive objective}
System 1 learns to predict structured next-step representations:
\[
\mathcal{L}_{AR} = -\sum_t \log p(y_t \mid y_{<t}, x)
\]

\subsection{RL objective}
RL updates policy/selection over candidate rules and inference actions:
\[
\nabla_\theta J(\theta) = \mathbb{E}\left[\sum_t \nabla_\theta \log \pi_\theta(a_t|s_t)\hat{A}_t\right]
\]
where rewards include delayed task outcomes and intermediate reasoning utility.

\subsection{Semantic-AR objective}
Semantic-AR adds reward shaping from semantic consistency:
\[
R_t = \alpha R^{task}_t + \beta R^{semantic}_t + \gamma R^{symbolic}_t
\]
with tunable weights for capability-vs-stability tradeoffs.

\section{Scoring and Rule Utility Philosophy}
Rule quality is treated as \textbf{soft and relative}, not binary.
Candidate rules remain in a large pool and accrue utility over trajectories.
Utility is updated from delayed outcomes and intermediate contributions.

\subsection{Candidate score decomposition (example)}
\[
\texttt{score}(r)=
w_c C(r)+w_s S(r)+w_u U(r)+w_t T(r)-w_d D(r)
\]
where $C$ is compressivity/complexity signal, $S$ symbolic consistency, $U$ learned utility, $T$ teacher/heuristic priors, and $D$ duplication/instability penalty.

\section{Experimental Program}

\subsection{Primary research question}
Can LT+RL(+semantic-AR), with symbolic scaffolding, achieve better reasoning-per-dollar than small dense LLM baselines?

\subsection{Benchmark layers}
\begin{itemize}[leftmargin=1.5em]
  \item Common-sense/general reasoning layer (primary).
  \item Arithmetic transfer layer (secondary diagnostic).
\end{itemize}

\subsection{Efficiency-first metrics}
\begin{itemize}[leftmargin=1.5em]
  \item task score per GPU dollar,
  \item score per training hour,
  \item score per parameter,
  \item improvement slope per added compute.
\end{itemize}

\section{Human-in-the-Loop Trial Governance}
Paid GPU trials should use explicit checkpoints:
\begin{enumerate}[leftmargin=1.5em]
  \item Pre-run budget cap confirmation.
  \item Post-first-run quality/cost review.
  \item Post-pilot go/no-go decision.
  \item Mid-batch stop/continue/scale decision.
  \item Post-batch budget forecast and next-phase decision.
\end{enumerate}

\section{Open Technical Questions}
\begin{itemize}[leftmargin=1.5em]
  \item What semantic-AR shaping terms provide stable gains without reward hacking?
  \item What rule-shape library gives best tradeoff between expressivity and GPU efficiency?
  \item How should delayed credit be distributed across chained WMEs and mixed-source rules?
  \item Which benchmark mix best reflects AGI-relevant progress under tight budgets?
\end{itemize}

\section{Planned Figure/Table Placeholders}
\begin{itemize}[leftmargin=1.5em]
  \item Figure: Tri-system architecture and WM/LTM data flow.
  \item Table: Rule-shape template library and intended semantic role.
  \item Table: Trial matrix (arms, seeds, cost caps, checkpoint policy).
  \item Figure: Capability-vs-cost frontier across runs.
\end{itemize}

\section{Conclusion}
This draft frames the project around an interpretable tri-system architecture with soft utility-driven rule learning,
semantic grounding, and cost-aware experimentation.
It is intended to be incrementally revised as evidence accumulates.

\end{document}

